\subsection{阿廷环上的模}

命题 2. —— 设 A 是一个左阿廷环。对每个 A-模 M，下列性质等价：

\begin{enumerate}
\def\labelenumi{(\roman{enumi})}
\item
  M 是半单的。
\item
  M 无根。
\item
  M 被 A 的根 r 零化。
\end{enumerate}

我们知道 (i) 蕴含 (ii)（VIII, 第 153 页，命题 3）。另一方面，rM 含于 M 的根（VIII, 第 158 页，命题 6），故 (ii) 蕴含 (iii)。

假设 A-模 M 被 r 零化。我们可以把它视为环 A/r 上的模。而环 A/r 是半单的（VIII, 第 173 页，命题 1），且半单环上的每个模都是半单的（VIII, 第 138 页，命题 4）。于是，M 是环 A/r 上的半单模，更是环 A 上的半单模。故 (iii) 蕴含 (i)。

推论. —— 设 A 是一个左阿廷环。对每个 A-模 M，M 的根等于 rM。

M 的根 \(\Re(\mathrm{M})\) 包含 rM（VIII, 第 158 页，命题 6）。此外，A-模 M/rM 被 r 零化。由命题 2，它因而无根，这蕴含 \(\Re(\mathrm{M}) \subset \mathfrak{rM}\)（VIII, 第 152 页，推论 1, c)）。

命题 3. —— 设 A 与 B 是环，f 是从 A 到 B 的一个同态，使得 B 由 \(f(\mathrm{A})\) 与交换子 \(f(\mathrm{A})'\)（即 \(f(\mathrm{A})\) 在 B 中的交换子）的并生成。设 M 是一个左 B-模。假设或者环 A 是左阿廷的，或者由 M 通过限制标量得到的 A-模 \(f_{*}(\mathrm{M})\)（II, §1, No.~13, 第 221 页）是有限生成的。则我们有 \(\mathfrak{R}_{\mathrm{A}}(f_{*}(\mathrm{M})) \subset \mathfrak{R}_{\mathrm{B}}(\mathrm{M})\)。

设 S 是一个单 B-模。假设或者 A 是左阿廷的，或者 A-模 \(f_{*}(S)\) 是有限生成的。我们将证明 \(f_{*}(S)\) 无根。对每个 \(b \in f(A)'\)，位似 \(b_{S}\) 是 A-模 \(f_{*}(S)\) 的一个自同态，因而保持 \(\mathfrak{R}_{\mathrm{A}}(f_{*}(\mathrm{S}))\)（VIII, 第 152 页，命题 1）。由于 B 由 \(f(\mathrm{A}) \cup f(\mathrm{A}')\) 生成，根 \(\mathfrak{R}_{\mathrm{A}}(f_{*}(\mathrm{S}))\) 是 S 的一个 B-子模，故等于 0 或 S。若 \(f_{*}(S)\) 是一个有限生成 A-模，则我们有 \(\mathfrak{R}_{\mathrm{A}}(f_{*}(\mathrm{S})) \neq f_{*}(\mathrm{S})\)（VIII, 第 153 页，命题 2）。若环 A 是左阿廷的，则我们有 \(\mathfrak{R}_{\mathrm{A}}(f_{*}(\mathrm{S})) = f(\mathfrak{r})\mathrm{S}\)，其中 r 是 A 的根（VIII, 第 174 页，命题 2 的推论），而由于 r 是 A 的一个幂零双边理想（VIII, 第 173 页，命题 1），我们不能有 \(S = f(\mathfrak{r})S\)。因此在两种情形下我们都有 \(\mathfrak{R}_{\mathrm{A}}(f_{*}(\mathrm{S})) = 0\)。

设 u 是从 M 到一个单 B-模 S 的非零 B-线性映射。映射 u 是满射；因而，若 \(f_{*}(\mathrm{M})\) 是有限生成的，则 \(f_{*}(\mathrm{S})\) 亦然。在命题 3 的假设下，A-模 \(f_{*}(\mathrm{S})\) 无根，且 u 是从 \(f_{*}(\mathrm{M})\) 到 \(f_{*}(\mathrm{S})\) 的一个 A-线性映射。于是由 VIII, 第 152 页的命题 1，u 的核包含 \(\mathfrak{R}_{\mathrm{A}}(f_{*}(\mathrm{M}))\)。由于 u 是任意的，我们有 \(\mathfrak{R}_{\mathrm{A}}(f_{*}(\mathrm{M})) \subset \mathfrak{R}_{\mathrm{B}}(\mathrm{M})\)。

推论. —— 设 A 是一个交换环，B 是 A 上的一个代数。假设或者环 A 是阿廷的，或者 B 是一个有限生成 A-模。则我们有 \(\Re(\mathrm{A})\mathrm{B}\subset\Re(\mathrm{B})\)。

由命题 3，根 \(\mathfrak{R}_{\mathrm{A}}(\mathrm{B}_{s})\) 含于 \(\mathfrak{R}_{\mathrm{B}}(\mathrm{B}_{s})\)，后者正是环 B 的根。此外，由 VIII, 第 158 页的命题 6，我们有 \(\mathfrak{R}(\mathrm{A})\mathrm{B}\subset\mathfrak{R}_{\mathrm{A}}(\mathrm{B}_{s})\)。推论得证。

